\section{Filtered maps and their successive obstructions}\label{s:chapter}

The common Jordan--H\"older theorem is already complete. A reader following
only that theorem may pass to Section~\ref{k:section}; the purpose here
is to extend the page construction to arbitrary filtered module maps.
We now develop the module page calculus beyond composition-series input.
Its simplest acyclic example will recover the preceding comparison.
We can put the two copies of an object into the exact complex
\[
0\longrightarrow M\xrightarrow{\id}M\longrightarrow0,
\]
give its two terms different filtrations, and calculate what happens to
their layers. The complex has no homology: the identity has zero kernel
and zero cokernel. Its initial layers need not cancel immediately,
because they occur at different filtration levels. A spectral sequence
records the successive cancellations.

When the layers are simple, every nonzero cancellation map is an
isomorphism. Since nothing survives, these maps pair every factor on one
side with exactly one factor on the other. This will prove the module
Jordan--H\"older theorem without assuming composition length, composition
multiplicities, or their invariance.

We shall construct all the pages needed for this argument. The setting is
left modules over a fixed ring with identity. Only submodules, quotient
modules, and module homomorphisms are used. The argument does not invoke
the earlier categorical Jordan--H\"older theorem. The preceding chapter
has already supplied the group and operator version of the particular
simple-factor calculation. Here additive module structure permits
arbitrary filtered maps, nonsimple layers, and nonzero limiting homology.

\subsection{A map seen one layer at a time}

Let $f:X\to Y$ be a module homomorphism. Give both modules finite
increasing filtrations indexed by the integers. Thus
\[
F_pX\subseteq F_{p+1}X,\qquad F_pY\subseteq F_{p+1}Y,
\]
the filtrations are zero for sufficiently negative $p$, and they equal
the whole modules for sufficiently positive $p$. Assume that $f$ preserves
them:
\[
f(F_pX)\subseteq F_pY\quad\text{for every }p.
\]
One may regard $X\xrightarrow{f}Y$ as a chain complex in degrees $1$ and
$0$. Its homology is
\[
H_1=\Ker f,\qquad H_0=Y/f(X).
\]
There are no other terms or differentials, so the condition that two
successive differentials compose to zero is automatic.

At the first level of resolution, keep only the associated layers:
\[
X_0(p)=F_pX/F_{p-1}X,\qquad
Y_0(p)=F_pY/F_{p-1}Y.
\]
Filtration preservation gives a well-defined module map
\[
d_{0,p}:X_0(p)\longrightarrow Y_0(p),\qquad [x]\longmapsto[f(x)].
\]
If $x$ is changed by an element of $F_{p-1}X$, its image changes by an
element of $F_{p-1}Y$. This is the entire well-definedness check.

Taking the kernel and cokernel of these maps retains the source classes
whose images have fallen to a lower layer, and the target classes not yet
accounted for by images. We then ask whether the retained source classes
reach a target one layer lower, two layers lower, and so on. The page
number measures how far down the filtration the map is being examined.

\subsection{The pages and their maps}

For $r\geq1$ and integers $p,q$, define
\begin{align}
X_r(p)
&=
\frac{F_pX\cap f^{-1}(F_{p-r}Y)}
     {F_{p-1}X\cap f^{-1}(F_{p-r}Y)},
\label{s:source-page}\\[3pt]
Y_r(q)
&=
\frac{F_qY}
     {F_{q-1}Y+\bigl(F_qY\cap f(F_{q+r-1}X)\bigr)}.
\label{s:target-page}
\end{align}
All the displayed sets are submodules. In the first formula the
denominator lies in the numerator because $F_{p-1}X\subseteq F_pX$;
in the second both summands of the denominator lie in $F_qY$.

The numerator of $X_r(p)$ consists of representatives in source layer $p$
whose image has reached target level $p-r$ or lower. Its denominator
identifies representatives that were already in a lower source layer
and satisfy the same condition. On the target side, $Y_r(q)$ discards
lower target layers and images already detected from source levels up
to $q+r-1$.

More precisely, $[x]\mapsto x+F_{p-1}X$ embeds $X_r(p)$ into
$F_pX/F_{p-1}X$: the map on its numerator has exactly the displayed
denominator as kernel.
The target $Y_r(q)$ is a quotient of $F_qY/F_{q-1}Y$. Thus a surviving
source entry is retained as a submodule of its original layer, while a
surviving target entry is retained as a quotient of its original layer.

There is a module map
\begin{equation}\label{s:page-differential}
d_{r,p}:X_r(p)\longrightarrow Y_r(p-r),\qquad[x]\longmapsto[f(x)].
\end{equation}
To check it, put $q=p-r$. A numerator representative has image in $F_qY$.
Changing it by
$z\in F_{p-1}X\cap f^{-1}(F_qY)$ changes the image by
\[
f(z)\in F_qY\cap f(F_{p-1}X)
       =F_qY\cap f(F_{q+r-1}X),
\]
which is in the target denominator. The formula is therefore independent
of representatives and is linear.

For each page, the arrows go from the source row to the target row:
\[
\begin{tikzcd}[column sep=large,row sep=large]
X_r(p) \arrow[d,"d_{r,p}"] & \text{source degree }1\\
Y_r(p-r) & \text{target degree }0.
\end{tikzcd}
\]
There are no arrows out of the target row or into the source row. Hence
the square of the differential is zero.

\begin{proposition}[The next page is homology]\label{s:transition}
There are canonical isomorphisms
\begin{align*}
X_{r+1}(p)&\cong\Ker\bigl(d_{r,p}:X_r(p)\to Y_r(p-r)\bigr),\\
Y_{r+1}(q)&\cong
Y_r(q)/\im\bigl(d_{r,q+r}:X_r(q+r)\to Y_r(q)\bigr)
\end{align*}
for every $r\geq0$.
\end{proposition}

\begin{proof}
For $r=0$, the kernel of $d_{0,p}$ consists of the classes represented by
$x\in F_pX$ with $f(x)\in F_{p-1}Y$. Since
$f(F_{p-1}X)\subseteq F_{p-1}Y$, its denominator is $F_{p-1}X$.
This is exactly $X_1(p)$. Quotienting $Y_0(q)$ by the image of $d_{0,q}$
gives
\[
F_qY/\bigl(F_{q-1}Y+f(F_qX)\bigr)=Y_1(q).
\]

Now take $r\geq1$ and put $q=p-r$. The numerator for $X_{r+1}(p)$ is
contained in the numerator for $X_r(p)$. Its intersection with the old
denominator is
\[
\bigl(F_pX\cap f^{-1}(F_{q-1}Y)\bigr)
\cap\bigl(F_{p-1}X\cap f^{-1}(F_qY)\bigr)
=F_{p-1}X\cap f^{-1}(F_{q-1}Y).
\]
That is precisely the new denominator. Thus inclusion of representatives
induces an injection $X_{r+1}(p)\hookrightarrow X_r(p)$.
Every class in its image maps to zero under $d_{r,p}$, because its image
under $f$ lies in $F_{q-1}Y$.

Conversely, suppose $[x]\in X_r(p)$ maps to zero. Then
\[
f(x)\in F_{q-1}Y+\bigl(F_qY\cap f(F_{p-1}X)\bigr).
\]
Choose $z\in F_{p-1}X$ with $f(z)\in F_qY$ such that
$f(x)-f(z)\in F_{q-1}Y$. The element $z$ is in the denominator of
$X_r(p)$, so $x-z$ represents the same class there. But $x-z$ is in the
numerator of $X_{r+1}(p)$. This proves that the injection identifies
$X_{r+1}(p)$ with the kernel.
The choice of $z$ proves that a suitable representative exists; the
canonical map itself is induced by inclusion and does not depend on that
choice.

On the target side, the denominators increase with $r$, giving a natural
surjection $Y_r(q)\twoheadrightarrow Y_{r+1}(q)$. The additional elements
killed are the classes of
\[
F_qY\cap f(F_{q+r}X).
\]
Every such element is $f(x)$ for an $x\in F_{q+r}X$ whose image lies in
$F_qY$, so it comes from $X_r(q+r)$. Conversely, every image of
$d_{r,q+r}$ has this form. The kernel of the surjection is therefore
exactly that image, proving the asserted cokernel formula.
\end{proof}

This calculation is a spectral sequence, in the ordinary algebraic
sense. To see the conventional indexing, set
\[
E_r^{p,q}=
\begin{cases}
X_r(p),&p+q=1,\\
Y_r(p),&p+q=0,\\
0,&\text{otherwise}.
\end{cases}
\]
Then $d_r$ has bidegree $(-r,r-1)$, the convention for an increasing
filtration on a chain complex, and Proposition~\ref{s:transition} says
exactly
\[
E_{r+1}=H(E_r,d_r).
\]
The explicit formulas establish this relation; no general spectral
sequence existence theorem is being assumed.
The support lies on the two total-degree diagonals $p+q=1$ and $p+q=0$.
The index $q$ may be negative; no first-quadrant assumption is needed.
The finite filtration support will supply the convergence bound.

The construction also respects maps. Suppose
\[
\begin{tikzcd}
X\arrow[r,"f"]\arrow[d,"a"']&Y\arrow[d,"b"]\\
X'\arrow[r,"f'"']&Y'
\end{tikzcd}
\]
commutes and $a,b$ preserve filtrations. The formula $[x]\mapsto[a(x)]$
defines $X_r(p)\to X'_r(p)$: the equality $f'a=bf$ carries both numerator
and denominator into their primed versions. Similarly $[y]\mapsto[b(y)]$
defines $Y_r(q)\to Y'_r(q)$, since
$b(f(F_jX))=f'(a(F_jX))\subseteq f'(F_jX')$. These maps commute with
$d_r$ by the same equality $f'a=bf$. The source inclusions and target
quotient maps defining the next page commute with $a,b$ as well, so
the identifications with homology are natural. Thus every page and differential
belongs to the given filtered map, rather than to choices of bases.

\subsection{What eventually survives}

\begin{proposition}[Convergence in this two-term setting]\label{s:convergence}
For sufficiently large $r$, the pages stabilize to
\begin{align*}
X_r(p)&=
\frac{F_pX\cap\Ker f}{F_{p-1}X\cap\Ker f},\\
Y_r(q)&=
\frac{F_qY}{F_{q-1}Y+(F_qY\cap f(X))}.
\end{align*}
The first is the associated layer of the induced filtration on $\Ker f$;
the second is canonically the associated layer of the quotient filtration
on $Y/f(X)$. In particular, if $f$ is an isomorphism, every page entry
is eventually zero.
\end{proposition}

\begin{proof}
Fix $p$. For sufficiently large $r$, $F_{p-r}Y=0$, so the two inverse
images in \eqref{s:source-page} become $\Ker f$. Fix $q$. For sufficiently
large $r$, $F_{q+r-1}X=X$, giving the second formula from
\eqref{s:target-page}.

Let $\pi:Y\to Y/f(X)$ be the quotient map. The quotient filtration is
$F_q(Y/f(X))=\pi(F_qY)$. The map
\[
F_qY\longrightarrow\pi(F_qY)/\pi(F_{q-1}Y),\qquad
y\longmapsto\pi(y)+\pi(F_{q-1}Y)
\]
is surjective. An element $y\in F_qY$ is in its kernel exactly when
$y=y_0+f(x)$ for $y_0\in F_{q-1}Y$. Then $f(x)=y-y_0\in F_qY$,
so the kernel is $F_{q-1}Y+(F_qY\cap f(X))$, as asserted.

Only finitely many filtration columns are nonzero on page $0$.
Proposition~\ref{s:transition} shows that a zero entry stays zero.
We may therefore take a single sufficiently large $r$ for all relevant
columns. If $f$ is an isomorphism, its kernel and cokernel are zero,
so all the stabilized entries are zero.
\end{proof}

Notice what the proof used: bounded filtrations, kernels, images, and
explicit quotient maps. It did not count dimensions or composition
factors. This is the point that makes a noncircular application to
Jordan--H\"older possible.

\subsection{Two composition series inside an acyclic complex}

Suppose that a module $M$ has two composition series
\[
0=A_0\subsetneq A_1\subsetneq\cdots\subsetneq A_m=M,
\qquad
0=B_0\subsetneq B_1\subsetneq\cdots\subsetneq B_n=M.
\]
Thus each $A_i/A_{i-1}$ and each $B_j/B_{j-1}$ is a nonzero simple
module. The numbers $m$ and $n$ are not assumed equal.

Extend the notation by $A_i=0$ for $i\leq0$, $A_i=M$ for $i\geq m$,
and similarly for $B$. Take $X=Y=M$ and $f=\id_M$, with
\begin{equation}\label{s:staggered-filtrations}
F_pX=A_{p-n},\qquad F_pY=B_p.
\end{equation}
The entire $B$ filtration occurs before the nonzero stages of the
$A$ filtration. This makes the identity filtration preserving:
if $p\leq n$, then $F_pX=0$; if $p>n$, then $F_pY=M$.
We have a filtered complex even though the inverse identity need not
preserve these two filtrations.

Its nonzero page-$0$ entries are precisely
\[
X_0(n+i)=A_i/A_{i-1}\quad(1\leq i\leq m),
\qquad
Y_0(j)=B_j/B_{j-1}\quad(1\leq j\leq n).
\]
Their column ranges are disjoint, so $d_0=0$ and $E_1=E_0$.
The positions can be summarized as follows:
\[
\begin{array}{c|c|c}
\text{row}&\text{nonzero columns}&\text{entry}\\ \hline
\text{source}&p=n+i,\quad1\leq i\leq m&A_i/A_{i-1}\\
\text{target}&p=j,\quad1\leq j\leq n&B_j/B_{j-1}
\end{array}
\]
Each $d_r$ travels from a source column to the target column $r$ places
to its left.

\begin{theorem}[Module Jordan--H\"older by spectral cancellation]
\label{s:jordan-holder}
The two composition series have equal length, $m=n$, and their factors
can be paired by module isomorphisms. The nonzero differentials in the
spectral sequence of \eqref{s:staggered-filtrations} supply such a pairing.
\end{theorem}

\begin{proof}
Recall first why a nonzero map $u:S\to T$ between simple modules is an
isomorphism. Its kernel is a submodule of $S$, and cannot be all of $S$
when $u\neq0$, so it is zero. Its image is a nonzero submodule of $T$,
so it is $T$. The first isomorphism theorem now gives the assertion.

Initially every nonzero page entry is simple. Inductively suppose this
is true on page $r$. A differential is either zero, leaving its source
and target unchanged on the next page, or is an isomorphism between two
simple entries, making both zero on the next page. There is only one
possible outgoing differential from each source position and one
possible incoming differential to each target position on a fixed page.
There are no incoming differentials to source positions and no outgoing
differentials from target positions. Proposition~\ref{s:transition}
therefore proves the induction.

Keep each original factor as a label on its position until that position
vanishes. A zero differential carries the label forward by the canonical
kernel or quotient identification. A nonzero differential pairs its two
labels by an isomorphism and removes them. Removed positions stay zero,
so no label is used twice. Proposition~\ref{s:convergence} applies to
$f=\id_M$ and says that eventually every position is zero. Consequently
every original label is paired exactly once, one $A$ factor with one
$B$ factor. This is a bijection of their finite index sets, proving
$m=n$ and the required isomorphisms.

If $M=0$, both composition series have no factors and the conclusion
holds directly.
\end{proof}

For $M\ne0$, termination also has an explicit bound: $E_{m+n}=0$.
Indeed, in a source column $p=n+i$, formula~\eqref{s:source-page}
then uses $F_{p-(m+n)}Y=B_{i-m}=0$. In a target column $j\geq1$,
formula~\eqref{s:target-page} uses
$F_{j+(m+n)-1}X=A_{j+m-1}=M$, so its denominator is all of $B_j$.
These bounds use the lengths of the two given lists, without assuming
that those lengths agree or are intrinsic to $M$.

This proof obtains an actual pairing. An equality of alternating sums
would give less information unless one had already established the
independence of simple classes. No such independence, or Grothendieck
group calculation, is needed here.

\subsection{A calculation with nonzero differentials}

Let $M=\kk e_1\oplus\kk e_2$, with
\[
A_1=\kk e_1,\qquad B_1=\kk e_2,\qquad A_2=B_2=M.
\]
Take $n=2$ in \eqref{s:staggered-filtrations}. Thus the source enters at
columns $3,4$, while the target enters at columns $1,2$:
\[
\begin{array}{c|cccc}
&1&2&3&4\\ \hline
X_1&0&0&\kk e_1&M/\kk e_1\\
Y_1&\kk e_2&M/\kk e_2&0&0.
\end{array}
\]
The only potentially nonzero $d_1$ is
\[
d_{1,3}:\kk e_1\longrightarrow M/\kk e_2,
\qquad ae_1\longmapsto ae_1+\kk e_2.
\]
It is an isomorphism: every target class has such a representative, and
the only vector common to the two coordinate lines is zero. Those two
positions disappear. Page $2$ therefore has just
\[
X_2(4)=M/\kk e_1,\qquad Y_2(1)=\kk e_2.
\]
There is no target in column $2$, so $d_{2,4}=0$; there is no source
in column $3$, so no $d_2$ reaches column $1$. Page $3$ has the same two
surviving factors, but its source formula gives a useful representative:
\[
X_3(4)=\frac{M\cap\kk e_2}{\kk e_1\cap\kk e_2}
       =\kk e_2,
\qquad
Y_3(1)=\frac{\kk e_2}{\kk e_2\cap\kk e_1}=\kk e_2.
\]
The identification of this $X_3(4)$ with $X_2(4)$ sends
$ae_2$ to $ae_2+\kk e_1$. Now
\[
d_{3,4}:\kk e_2\xrightarrow{\id}\kk e_2
\]
removes the remaining pair, and $E_4=0$.
In particular, the zero differential $d_2$ did not mean that the
calculation had stabilized: $d_3$ was still nonzero. Equality of two
successive pages does not by itself imply that every later differential
vanishes.

The cancellations occur at different distances:
\[
\begin{tikzcd}[column sep=large,row sep=large]
\kk e_1\arrow[r,"d_{1,3}","\sim"']&M/\kk e_2
&\text{columns }3\to2\\
\kk e_2\arrow[r,"d_{3,4}","\sim"']&\kk e_2
&\text{columns }4\to1.
\end{tikzcd}
\]
This last picture records the two nonzero arrows on their respective
pages; it is not a single page. The source label $\kk e_2$ in the second row
is the page-$3$ representative just calculated.

For comparison, if $A_1=B_1=\kk e_1$, then $d_1=0$ and the two nonzero
maps occur on page $2$:
\[
d_{2,3}:\kk e_1\iso\kk e_1,
\qquad
d_{2,4}:M/\kk e_1\iso M/\kk e_1.
\]
The page of cancellation records relative filtration position; it is not
an invariant of the unfiltered module. One may also give $M$ the operator
$Te_1=0$, $Te_2=e_2$ and regard it as a $\kk[t]$-module. The same crossed
filtration calculation then pairs the simple type with eigenvalue $0$
on page $1$, and the different simple type with eigenvalue $1$ on page
$3$. The argument matches module types, as well as the number of terms.

\subsection{The butterfly is visible inside the differentials}

The construction has a precise relation to the earlier refinement
rectangle. Take any two finite filtrations $A$ and $B$ of a module $M$;
their factors need not be simple. Use the same shifted identity complex.
For $1\leq i\leq m$ and $1\leq j\leq n$, put
\[
r=n+i-j,\qquad p=n+i.
\]
Then $r\geq1$, and the explicit pages become
\begin{align*}
X_r(n+i)&=\frac{A_i\cap B_j}{A_{i-1}\cap B_j},\\
Y_r(j)&=\frac{B_j}{B_{j-1}+(A_{i-1}\cap B_j)}.
\end{align*}
The differential is induced by the inclusion $A_i\cap B_j\hookrightarrow
B_j$. Its image is canonically
\begin{equation}\label{s:differential-cell}
\im d_{n+i-j,n+i}\cong
Q_{ij}:=
\frac{A_i\cap B_j}
{(A_{i-1}\cap B_j)+(A_i\cap B_{j-1})}.
\end{equation}
Here is the kernel check. If $x\in A_i\cap B_j$ has zero target class,
write $x=b+z$ with $b\in B_{j-1}$ and $z\in A_{i-1}\cap B_j$.
Since both $x$ and $z$ lie in $A_i$, so does $b=x-z$; hence
$b\in A_i\cap B_{j-1}$. The converse is immediate. Thus the kernel of
the map from $A_i\cap B_j$ to its image is exactly the denominator
in \eqref{s:differential-cell}, and the first isomorphism theorem applies.

The same $Q_{ij}$ is what appears on taking associated layers in either
order. To state this precisely, filter $A_i/A_{i-1}$ by
\[
\overline B_j=\frac{A_{i-1}+(A_i\cap B_j)}{A_{i-1}}.
\]
The map $A_i\cap B_j\to\overline B_j/\overline B_{j-1}$ is surjective.
Its kernel is
$(A_{i-1}\cap B_j)+(A_i\cap B_{j-1})$: if
$x=a+y$ with $a\in A_{i-1}$ and $y\in A_i\cap B_{j-1}$, then
$a=x-y\in B_j$; the reverse inclusion is immediate. Consequently
\[
\operatorname{gr}^{B}_j\operatorname{gr}^{A}_i M
\cong Q_{ij}
\cong\operatorname{gr}^{A}_i\operatorname{gr}^{B}_j M.
\]
The second isomorphism follows by the same calculation after interchanging
$A$ and $B$. These are isomorphisms of individual indexed subquotients;
no simultaneous choice of a basis or a splitting of the filtrations has
been made.

Thus the static comparison rectangle and the spectral calculation share
the same subquotients. The rectangle displays all $Q_{ij}$ together.
The spectral sequence encounters them in order of the distance
$r=n+i-j$, as images of actual differentials. The kernel and quotient
calculations proved above explain the connection. Calling the rectangle
alone a spectral sequence would omit the page transitions and the
differentials that we have now supplied.

\subsection{What this adds, and what it does not assert}

The initial framework is a filtered map. Its successive pages separate
what the map already accounts for from what remains. The stable pages
retain the layers of its kernel and cokernel. When the map is an identity,
these stable pages are zero; when the initial layers are simple, the
necessary cancellations become isomorphisms. Jordan--H\"older emerges
from those two facts together.

A single filtration on a module concentrated in degree zero has no
nonzero differentials. Its spectral sequence leaves its own associated
graded object in place and does not by itself compare a second
filtration. The two-term identity complex is what creates a comparison
with nonzero differentials. Likewise, the bare statement that an
alternating sum vanishes would not supply our pairing proof.

With finite filtrations, the stabilized pages describe associated
layers of homology. They do not, without further data, specify how those
layers are attached. That limitation returns us to the questions about
extensions and reconstruction in the following chapters. Spectral
sequences organize successive information; they do not make the extension
problem disappear.

For the general theory, the Stacks Project constructs filtered-complex
spectral sequences and proves finite-filtration convergence in
\href{https://stacks.math.columbia.edu/tag/012K}{Section 12.24,
especially Lemmas 12.24.2 and 12.24.11}. Its convention is cohomological
with decreasing filtrations, whereas ours is homological with increasing
filtrations. The two-filtration interchange is explicitly identified
with Zassenhaus's lemma in Deligne's
\href{https://www.numdam.org/item/10.1007/BF02684692.pdf}{\emph{Th\'eorie
de Hodge II}, \S1.2.1}. These references place the constructions in their
classical setting; the proof in this chapter is given directly, and no
originality claim is made for this presentation.
